\chapter{Experimental Results \& Comprehensive Ablations}
\label{ch:results}

\draftnote{\textbf{This chapter is deliberately held as a data contract, not prose.} Two of
  its inputs are still in motion --- the fine-tuned-arm repeat sweep
  (Section~\ref{sec:res-encoder}) and the ADNI external validation
  (Section~\ref{sec:external-validation}) --- and the headline framing of the thesis is to
  be decided \emph{after} they land rather than before. Every section below therefore
  declares the table or figure it will contain, the artifact that will supply it, and the
  condition under which it may be filled. Nothing is written here that the artifacts do not
  yet support. Chapters~\ref{ch:foundations}, \ref{ch:data}, \ref{ch:methodology}
  and~\ref{ch:sota-repair} are complete and can be read independently of this one.}

%==============================================================================
\section{Reporting Rules for This Chapter}
\label{sec:reporting-rules}

These rules are fixed in advance so that no table in this chapter can be assembled in a way
that flatters a result.

\begin{enumerate}
  \item \textbf{Every table carries the floors.} The metadata-only baseline
    (CV AUC $0.6157 \pm 0.0653$, test AUC $0.4929$) and the raw-FC representation floor
    appear as rows in every model comparison, not in a separate section.
  \item \textbf{Every model number is a distribution, not a point.} Cross-validation results
    are reported over folds; multi-seed results over seeds; and where an arm is known to be
    non-deterministic, over \emph{repeats nested in seeds}.
  \item \textbf{Sample standard deviation throughout} ($\mathtt{ddof}=1$), so that error
    bars across documents are comparable.
  \item \textbf{No $p$-value is reported where its independence assumption fails.}
    Specifically: no paired test over cross-validation folds within a seed (folds share
    training data), and no paired test over seeds for an arm whose within-seed variance is
    unmeasured or dominant (Section~\ref{sec:btgt-stats}).
  \item \textbf{Every reported number globs all runs for its experiment identifier.} No
    single run path is hard-coded, since that is precisely how an unrepresentative run
    previously became ``the'' result (Section~\ref{sec:btgt-stats}).
  \item \textbf{Test-set size is printed beside every test metric}, because at
    $N_{\text{test}} = 34$ (or 25 in the matched window) one subject is worth
    $\sim 0.03$ AUC.
\end{enumerate}

%==============================================================================
\section{Longitudinal Conversion Prediction Benchmark}
\label{sec:res-benchmark}

\draftnote{\textbf{Table 6.1 --- main DELCODE benchmark.} Rows: metadata floor, raw-FC
  floor, GEP, GEC, GE-GRU, GE-LSTM. Columns: CV AUC (mean $\pm$ SD over 5 folds), test AUC,
  sensitivity, specificity, balanced accuracy, trainable parameters. Source:
  \texttt{CLASSIFIER/outputs/RESULTS.csv}, columns \texttt{cv.val\_auc\_mean},
  \texttt{cv.val\_auc\_std}, \texttt{metric.test\_*}. \emph{Fill condition:} immediately ---
  all runs exist.}

\risk{The static GEC classifier currently outscores every recurrent model on this cohort,
and the thesis must confront that rather than table it.}{%
\texttt{gec-trajectory-whole-brain} reaches test AUC $0.964$ against GE-LSTM's $0.921$ and
GE-GRU's $0.832$ on the full cohort. If a per-visit classifier with an explicit visit-count
mask beats the recurrent models, then either (a) the recurrence is not contributing, or
(b) GEC is exploiting the length shortcut that its mask makes available and the recurrent
models do not receive (Section~\ref{sec:static-baselines}). The within-subject slope
diagnostics of Section~\ref{sec:res-horizon} are what distinguishes these, and this section
must state the answer explicitly. Quietly reporting the recurrent models as the headline
while GEC sits higher in the same table is the single most likely thing to be challenged in
a defence.}

%==============================================================================
\section{Comparative Evaluation: GE-LSTM vs.\ Brain-TokenGT}
\label{sec:res-comparison}

\draftnote{\textbf{Table 6.2 + Figure 6.1 --- matched-window head-to-head.} Both models
  restricted to $2 \le T \le 3$ ($N_{\text{CV}} = 95$, $N_{\text{test}} = 25$), byte-matched
  subject identifiers and fold membership. GE-LSTM: 4 seeds, one run each. Brain-TokenGT:
  \textbf{all 19 runs}, reported as repeats nested in seeds per
  Table~\ref{tab:btgt-19runs}. Report as two distributions plus an effect size. \emph{No
  $p$-value} --- see Section~\ref{sec:btgt-stats} for why both previously computed tests are
  withdrawn. Figures already generated under
  \texttt{DOCS/per\_section/results/figures/}: ROC/PR curves, confusion matrices,
  calibration, parameter-efficiency frontier. \emph{Fill condition:} after the comparison
  notebook is re-executed with the glob-all-runs fix (Phase B).}

\draftnote{\textbf{Table 6.3 --- the reproducibility contrast.} GE-LSTM \texttt{none} arm:
  two runs at seed 42, byte-identical. Brain-TokenGT stabilised: 7 runs at seed 42, range
  $0.3571$--$0.7078$. This is a result in its own right and should be presented as one, not
  as a caveat inside Table 6.2.}

%==============================================================================
\section{Encoder Ablation: Does the Graph Encoder Earn Its Place?}
\label{sec:res-encoder}

\draftnote{\textbf{Table 6.4 --- the four encoder arms} of
  Table~\ref{tab:encoder-arms}, across seeds 42--45, against both floors. Source:
  \texttt{recon-ablation-gelstm-\{none,random,pretrained-frozen,pretrained-finetuned\}}
  and their \texttt{-seed4\{3,4,5\}} repeats. \emph{Fill condition:} the
  \texttt{pretrained\_finetuned} arm is confirmed non-deterministic (folds diverged on
  re-run at seed 43), so its row must report repeats nested in seeds, not one run per
  seed. The repeat sweep is running; this table is filled when it completes.}

\draftnote{\textbf{Read the result against the pre-registered interpretation table}
  in \texttt{DOCS/reconstruction-value-ablation.md} --- reproduce that table here verbatim
  and mark which row the observed pattern matches, before stating any conclusion. That
  ordering is the point of pre-registration and must be visible to the reader.}

\risk{The current fine-tuning conclusion is withdrawn, not merely caveated.}{%
The earlier reading --- that fine-tuning ``destabilises optimisation and severely
overfits'' --- rests on an arm in which the pretrained GATv2 encoder and the freshly
initialised classifier head were optimised at a \emph{single shared learning rate}. That is
the textbook way to destroy pretrained features, and it is precisely the treatment
Brain-TokenGT's own stabilisation applied to \emph{its} newly unfrozen parameters and this
arm did not. The arm therefore measures naive fine-tuning, not fine-tuning. Until either a
learning-rate-scaled arm is run or the claim is restated, no conclusion about fine-tuning
belongs in this chapter.}

%==============================================================================
\section{Trajectory Ablations: Order, Timing, and Horizon}
\label{sec:res-horizon}

\draftnote{\textbf{Table 6.5 --- eval-time trajectory ablations} against a fixed trained
  checkpoint: baseline; shuffled visit order ($5$ repeats); $\Delta t$ zeroed; first-2-visit
  window; metadata floor. Source: \texttt{sanity-gelstm-ablations}. These results
  \emph{exist} and are partially reported already in Section~\ref{sec:delta-t-finding};
  this section presents the full table.}

\draftnote{\textbf{Table 6.6 + Figure 6.2 --- AUC versus visit horizon, both
  constructions.} Variable-cohort and fixed-cohort curves side by side with per-$N$
  converter/non-converter composition printed beside each point
  (Section~\ref{sec:gap-protocol}). The fixed-cohort curve is the claim-bearing one.
  Source: \texttt{common/early\_detection.py} and \texttt{common/visit\_confound.py} via
  \texttt{sanity-visit-confound-*}.}

\draftnote{\textbf{Table 6.7 --- the follow-up-duration confound.} Visit-count distribution
  by group (converters 3.43 vs.\ stable 2.74 scans, Mann--Whitney $p = 0.002$); between-subject
  Spearman correlations; and the decisive within-subject probability slopes. This is what
  answers the GEC exposure raised in Section~\ref{sec:res-benchmark}, and it must be
  reported for GEC, GE-LSTM and GE-GRU together.}

%==============================================================================
\section{External Validation on ADNI}
\label{sec:external-validation}

\draftnote{\textbf{Table 6.8 --- ADNI replication.} Two arms only --- \texttt{none} and
  \texttt{pretrained\_frozen} --- across 4 seeds, on the ADNI longitudinal cohort
  ($162$ subjects with $\ge 2$ sessions, $51$ converters). \emph{Fill condition:} the
  cohort-aware visit parser must be wired into the GE-LSTM dataset \emph{and} pass the
  byte-equality gate on DELCODE first (Section~\ref{sec:visit-semantics}); if the gate fails,
  this section is cut and moved to Chapter~\ref{ch:conclusion}.}

\draftnote{\textbf{This section carries two claims that DELCODE cannot support.} First,
  whether the encoder ablation result of Section~\ref{sec:res-encoder} is a property of the
  method or of the cohort. Second --- and only here --- whether the $\Delta t$ conditioning
  contributes anything at all, since ADNI's follow-up is genuinely irregular where
  DELCODE's is 90\% uniform (Section~\ref{sec:delta-t-finding}). Report the ADNI inter-visit
  interval distribution beside Table~\ref{tab:delcode-intervals} so the contrast is visible.}

%==============================================================================
\section{Latent Space Interpretability}
\label{sec:res-interpretability}

\draftnote{\textbf{Figure 6.3 --- UMAP projections} of the frozen GAAE and VGAE latent
  spaces, coloured by diagnosis and by visit. \textbf{Figure 6.4 --- Integrated Gradients
  attribution} (\texttt{captum}) over the frozen encoders, aggregated to Schaefer-200
  parcels and summarised by Yeo network. Source: existing run artifacts under
  \texttt{CLASSIFIER/outputs/explain-gaae/} and \texttt{explain-vgae-gcn/}.}

\risk{Interpretability of an encoder whose contribution is unproven needs framing, not
suppression.}{%
If the encoder ablation shows that removing the graph encoder does not hurt performance,
then attribution maps computed \emph{over that encoder} explain a component that may not be
carrying the prediction. The honest framing is that these maps characterise what the
pretrained representation encodes --- a legitimate descriptive question --- and not what
drives the classifier's decisions. Presenting them as biomarker discovery would overstate
them.}

%==============================================================================
\section{Ablations Not Run}
\label{sec:res-not-run}

\edge{Stating which planned ablations were not run is part of the result.}{%
The \textbf{parcellation ablation} (whole brain vs.\ DMN vs.\ DMN + limbic/hippocampus) was
planned and is \emph{not} reported: no such experiment exists in the registry, and
fabricating or hand-waving it would be worse than omitting it. It is carried to
Chapter~\ref{ch:conclusion} as future work. Likewise the imaging-based survival head
(Section~\ref{sec:prognoser-scope}). A reader can verify both claims against the experiment
registry, which is the point.}
